\documentclass[a4paper]{article}

\usepackage[fontset=fandol]{ctex}
\usepackage{amsmath, amssymb}
\usepackage{graphicx}
\usepackage[margin=2.45cm]{geometry}
\usepackage[most]{tcolorbox}
\usepackage{hyperref}
\usepackage{booktabs}
\usepackage{float}
\usepackage{array}
\usepackage{xcolor}
\usepackage{enumitem}
\usepackage{caption}

\setlength{\parindent}{2em}
\setlength{\parskip}{0.35em}
\setlength{\emergencystretch}{3em}
\linespread{1.06}
\sloppy
\raggedbottom
\vfuzz=4pt

\newcolumntype{L}[1]{>{\raggedright\arraybackslash}p{#1}}
\newcolumntype{C}[1]{>{\centering\arraybackslash}p{#1}}

\newcommand{\E}{\mathbb{E}}
\newcommand{\oldtheta}{\theta^{b}}
\newcommand{\newtheta}{\theta}
\newcommand{\loss}{L}
\newcommand{\grad}{g}

\hypersetup{
    colorlinks=true,
    linkcolor=blue!60!black,
    urlcolor=blue!60!black
}

\setlist[itemize]{leftmargin=2.2em,itemsep=0.22em,topsep=0.25em}
\setlist[enumerate]{leftmargin=2.4em,itemsep=0.22em,topsep=0.25em}
\setlist[description]{leftmargin=3.0em,itemsep=0.18em,topsep=0.25em}

\newtcolorbox{knowledgebox}[1]{
    enhanced, colback=blue!5!white, colframe=blue!75!black, colbacktitle=blue!75!black,
    coltitle=white, fonttitle=\bfseries, title={#1},
    attach boxed title to top left={yshift=-2mm, xshift=2mm},
    boxrule=1pt, sharp corners
}

\newtcolorbox{importantbox}[1]{
    enhanced, colback=yellow!10!white, colframe=yellow!80!black, colbacktitle=yellow!80!black,
    coltitle=black, fonttitle=\bfseries, title={#1}, sharp corners
}

\newtcolorbox{warningbox}[1]{
    enhanced, colback=red!5!white, colframe=red!75!black, colbacktitle=red!75!black,
    coltitle=white, fonttitle=\bfseries, title={#1}, sharp corners
}

\newcommand{\notetitle}{机器终身学习（二）：缓解灾难性遗忘}
\newcommand{\notesubtitle}{Selective Synaptic Plasticity, GEM, Neural Resource Allocation, Memory Replay, and Curriculum}
\newcommand{\noteauthors}{基于李宏毅老师《机器学习 2025》课程内容整理}
\newcommand{\notedate}{\today}
\newcommand{\videochannel}{李宏毅深度学习课堂（Bilibili）}
\newcommand{\videoduration}{约 36 分 39 秒}
\newcommand{\videourl}{https://www.bilibili.com/video/BV1TAtwzTE1S?p=54}

\newcommand{\framefig}[4]{%
\begin{figure}[H]
    \centering
    \includegraphics[width=#1\textwidth]{#2}
    \caption{#3\protect\footnotemark}
\end{figure}
\footnotetext{视频画面时间区间：#4。}
}

\begin{document}
\pagestyle{empty}

\begin{center}
    \vspace*{1.0cm}
    {\Huge\bfseries \notetitle\par}
    \vspace{0.35cm}
    {\LARGE \notesubtitle\par}
    \vspace{0.75cm}
    \includegraphics[width=0.62\textwidth]{video.jpg}\par
    \vspace{0.75cm}
    {\large \noteauthors\par}
    {\large \notedate\par}
    \vspace{0.75cm}
    \begin{tcolorbox}[width=0.86\textwidth,center,sharp corners,
                      colback=black!3!white,colframe=black!50!white]
        \centering
        \textbf{视频信息}\\[3pt]
        频道：\videochannel \\
        时长：\videoduration \\
        链接：\href{\videourl}{\videourl}
    \end{tcolorbox}
\end{center}

\newpage
\tableofcontents
\newpage
\pagestyle{plain}

% =====================================================================
\section{缓解遗忘的三类研究方向}
% =====================================================================

P53 说明了 catastrophic forgetting：模型依序学习任务时，后续任务会把旧任务能力冲掉。P54 开始讨论解决方向。课程先给出三类方法：
\begin{itemize}
    \item \textbf{Selective Synaptic Plasticity：} 选择性地保护旧任务重要参数，只让不重要参数更自由地变化。
    \item \textbf{Additional Neural Resource Allocation：} 为新任务分配额外网络资源，避免所有任务抢同一组参数。
    \item \textbf{Memory Replay：} 通过保留、生成或回放旧任务数据，让新任务训练时仍看见旧任务。
\end{itemize}

\framefig{0.84}{figures/fig01_research_directions_three_families.jpg}
{缓解 catastrophic forgetting 的三大方向：selective synaptic plasticity、additional neural resource allocation 与 memory replay；课程首先聚焦 regularization-based approach。}
{00:01:30--00:01:45}

这三类方法对应了遗忘问题的三个切入点。Selective synaptic plasticity 直接约束参数更新；additional resources 改变模型容量和任务分配方式；memory replay 则改变训练数据流，让模型在学新任务时仍接触旧任务信号。

\begin{importantbox}{三类方法的核心差异}
    正则化方法通常不保存大量旧数据，代价低但保护能力有限；资源分配方法能显式隔离任务，但模型会增长或可用参数减少；replay 方法最接近 multi-task training，但需要存储旧样本或训练生成模型。
\end{importantbox}

\subsection{本章小结}

P54 的方法地图围绕三个问题展开：哪些参数该被保护、是否给新任务新资源、以及如何让旧任务数据继续参与训练。后续每类方法都在 stability 与 plasticity 之间做取舍。

% =====================================================================
\section{为什么会发生 Catastrophic Forgetting}
% =====================================================================

课程用二维参数空间解释遗忘。设模型参数为 $(\theta_1,\theta_2)$，task 1 与 task 2 有不同 error surface，颜色越深表示 loss 越小。模型先从 $\theta^0$ 训练到 task 1 的好解 $\theta^b$；接着训练 task 2 时，梯度会把参数推向 task 2 的低 loss 区域 $\theta^\ast$。

\framefig{0.84}{figures/fig02_error_surface_forgetting.jpg}
{两个任务的 error surface 不同：从旧任务解 $\theta^b$ 往 task 2 低 loss 区移动时，参数可能离开 task 1 的低 loss 区，导致 task 1 被遗忘。}
{00:05:15--00:05:30}

遗忘的关键不在于 task 2 本身有多难，而在于 task 2 的好解和 task 1 的好解不一定重合。若优化器只看当前任务 loss，它会沿着降低 task 2 loss 的方向更新；这个方向可能同时让 task 1 loss 上升。越多任务共享同一组参数，这种冲突越容易出现。

从这个图也能看出几种可能策略：
\begin{itemize}
    \item 保护 task 1 重要方向，不让参数离开旧任务低 loss 区太远。
    \item 找一个对 task 2 有利、但不增加 task 1 loss 的更新方向。
    \item 给 task 2 新参数，减少它和 task 1 对旧参数的争抢。
    \item 在训练 task 2 时继续用 task 1 数据提醒模型。
\end{itemize}

\subsection{本章小结}

Catastrophic forgetting 可以看成参数空间中的任务冲突：当前任务梯度把参数推向新任务低 loss 区，却可能推离旧任务低 loss 区。P54 的方法本质上都在限制、重定向或补偿这种参数移动。

% =====================================================================
\section{Selective Synaptic Plasticity：保护重要参数}
% =====================================================================

Selective Synaptic Plasticity 的基本思想是：模型中并非所有参数对旧任务都同等重要。某些参数如果改变，旧任务性能会大幅下降；另一些参数即使改变，旧任务也不太受影响。因此，学习新任务时应更多修改不重要参数，尽量保护重要参数。

课程给出一个正则化形式：
\[
    L'(\theta)
    =
    L(\theta)
    +
    \lambda \sum_i b_i(\theta_i-\theta_i^b)^2.
\]
其中 $\theta^b$ 是从旧任务学到的参数，$b_i$ 是第 $i$ 个参数的 guard 或 importance，$\lambda$ 控制保护强度。若 $b_i$ 大，$\theta_i$ 离开 $\theta_i^b$ 会受到更大惩罚；若 $b_i$ 小，该参数更容易被新任务更新。

\framefig{0.84}{figures/fig03_selective_synaptic_plasticity_loss.jpg}
{Selective synaptic plasticity 的正则化目标：当前任务 loss 加上参数偏离旧任务参数的惩罚，每个参数的惩罚由重要性 $b_i$ 加权。}
{00:08:00--00:08:15}

这个式子可以看成从普通 L2 regularization 扩展而来。普通 L2 往往对所有参数同等惩罚；selective synaptic plasticity 则让每个参数有不同保护强度。真正困难的是估计 $b_i$：哪些参数对旧任务重要？重要到什么程度？

\subsection{保护强度的两端：遗忘与僵化}

如果 $b_i=0$，第 $i$ 个参数没有任何约束，学习新任务时可以随意改变。这提升 plasticity，但可能造成 catastrophic forgetting。若 $b_i=\infty$，参数永远不能离开旧值 $\theta_i^b$，旧任务被保护，但新任务可能学不动，这叫 intransigence。

\framefig{0.84}{figures/fig04_importance_guard_tradeoff.jpg}
{参数重要性 $b_i$ 的两端：$b_i=0$ 时没有约束，容易遗忘；$b_i=\infty$ 时参数不能改变，模型会过度僵化，难以学习新任务。}
{00:11:30--00:11:45}

因此 continual learning 的目标不是把所有旧参数冻结，而是估计“哪些方向必须靠近旧解，哪些方向可以让给新任务”。课程图中用两维参数说明：若 $b_1$ 小而 $b_2$ 大，模型可以主要改变 $\theta_1$，但尽量不改变 $\theta_2$，从而在 task 2 中寻找不破坏 task 1 的解。

\framefig{0.84}{figures/fig05_selective_plasticity_geometry.jpg}
{选择性保护的几何直觉：当 $b_1$ 小、$b_2$ 大时，可以沿 $\theta_1$ 方向调整模型，但不要大幅改变 $\theta_2$；这限制了新任务更新的可行方向。}
{00:14:45--00:15:00}

\subsection{EWC 与重要性估计}

Elastic Weight Consolidation（EWC）是代表性方法之一。它用 Fisher information 估计每个参数对旧任务 loss 的敏感程度：若某参数附近的 loss 曲率大，改变它会显著破坏旧任务，因此 $b_i$ 应该大。课程图中比较 EWC、普通 L2 和 SGD，在 permuted MNIST 的 task A、B、C 上，EWC 更能维持旧任务表现。

\framefig{0.84}{figures/fig06_ewc_l2_sgd_results.jpg}
{EWC、L2 与 SGD 的比较：在连续训练 task A、B、C 时，EWC 通过参数重要性保护旧任务，通常比普通 SGD 更不易遗忘；但过度保护会带来 intransigence。}
{00:18:30--00:18:45}

课程也列出一组相关方法：EWC、Synaptic Intelligence（SI）、Memory Aware Synapses（MAS）、RWalk、Sliced Cramer Preservation（SCP）。它们的共同点是：不直接保存大量旧数据，而是用某种方式估计旧任务对参数的约束，再在新任务训练中加入正则项。

\framefig{0.82}{figures/fig07_selective_plasticity_methods.jpg}
{Selective synaptic plasticity 的相关方法：EWC、SI、MAS、RWalk 与 SCP 等，核心都在于估计并保护旧任务重要参数。}
{00:22:15--00:22:30}

\begin{warningbox}{正则化保护无法凭空创造容量}
    如果新旧任务确实需要冲突很大的参数配置，单靠正则项只能在遗忘和僵化之间折中。保护太弱会忘，保护太强会学不动；保护参数也不能解决旧数据分布完全不可恢复的问题。
\end{warningbox}

\subsection{本章小结}

Selective synaptic plasticity 用重要性加权的正则项保护旧任务参数。它的优点是额外存储较少、实现相对直接；缺点是需要估计参数重要性，并且在任务冲突强时可能面临遗忘与 intransigence 的两难。

% =====================================================================
\section{Gradient Episodic Memory：重定向当前梯度}
% =====================================================================

GEM（Gradient Episodic Memory）从梯度角度处理遗忘。它保留一小部分旧任务数据，训练当前任务时也计算旧任务上的梯度。若当前任务梯度会让旧任务 loss 上升，就把当前梯度投影到一个不伤害旧任务的方向。

设当前任务负梯度为 $g$，旧任务 memory 上的负梯度为 $g^b$。GEM 希望找到一个尽量接近 $g$ 的更新方向 $g'$，同时满足
\[
    {g'}^\top g^b \ge 0.
\]
这个约束表示：新更新方向与旧任务下降方向不冲突，至少在一阶近似下不会增加旧任务 loss。更一般地，对每个旧任务 $k$ 都有约束 ${g'}^\top g_k \ge 0$。

\[
    \min_{g'} \|g'-g\|_2^2
    \quad
    \text{s.t.}\quad
    {g'}^\top g_k \ge 0,\ \forall k< t.
\]

\framefig{0.84}{figures/fig08_gem_gradient_projection.jpg}
{GEM 的几何图：红色虚线是当前任务负梯度，绿色是旧任务负梯度，红色实线 $g'$ 是投影后的更新方向；约束 $g'\cdot g^b\ge 0$ 需要旧任务数据来估计。}
{00:26:00--00:26:15}

GEM 和 EWC 的差别在于，EWC 用参数重要性近似旧任务约束，GEM 则直接使用旧任务样本估计“当前更新会不会伤害旧任务”。因此 GEM 更接近 replay 方法，但它不是简单把旧样本混进 batch，而是把旧样本用于约束梯度方向。

\begin{importantbox}{GEM 的关键代价}
    GEM 需要保存旧任务样本或代表性 memory。memory 越大，对旧任务约束越可靠，但存储和计算越高；memory 太小，则旧任务梯度可能不能代表完整旧分布。
\end{importantbox}

\subsection{本章小结}

GEM 把 continual learning 写成一个带约束的梯度更新问题：尽量沿当前任务下降方向走，但不能让旧任务 loss 上升。它比纯正则化更直接利用旧任务数据，也因此付出 memory 与额外梯度计算成本。

% =====================================================================
\section{Additional Neural Resource Allocation：给新任务分配资源}
% =====================================================================

第二大类方法不再只限制旧参数，而是为新任务分配额外神经资源。直觉是：如果所有任务都挤在同一组参数里，冲突不可避免；若新任务有自己的参数子空间，旧任务参数就不必被大幅修改。

\subsection{Progressive Neural Networks}

Progressive Neural Networks 为每个新任务增加一个新 column。旧任务 column 被冻结，避免遗忘；新任务 column 可以通过 lateral connections 读取旧任务表示，从而迁移旧知识。这样每个任务都有独立路径，又能利用过去学到的特征。

\framefig{0.84}{figures/fig09_progressive_neural_networks.jpg}
{Progressive Neural Networks：每个任务新增一列网络，旧列冻结，新列通过横向连接使用旧任务表示；这样避免遗忘，但模型规模随任务数增长。}
{00:28:30--00:28:45}

优点是遗忘少，因为旧任务参数不再被更新；缺点也明显：任务越多，模型越大。如果任务流很长，参数量和推理成本都会线性增长。它适合任务数量有限且任务边界明确的场景。

\subsection{PackNet、CPG 与参数复用}

PackNet 和 Compacting, Picking, and Growing（CPG）试图更节省地分配参数。PackNet 的思路是：先为一个任务训练网络，再剪枝释放一部分不重要参数；保留并冻结对旧任务重要的参数，空出来的参数给新任务使用。CPG 则在 compact、pick、grow 之间循环，必要时扩展模型。

\framefig{0.86}{figures/fig10_packnet_cpg_allocation.jpg}
{PackNet 与 CPG 的参数分配思路：通过 pruning、mask、free weights 与必要扩展，为不同任务保留或分配参数子空间。}
{00:30:30--00:30:45}

这类方法位于“完全冻结旧模型”和“所有任务共享全部参数”之间。它们让任务共享一部分网络，同时为关键任务保留专属参数。挑战在于任务顺序、剪枝比例、mask 管理和容量耗尽：如果早期任务占用太多参数，后续任务可能没有足够空间。

\subsection{本章小结}

Additional neural resource allocation 用结构隔离缓解遗忘。Progressive networks 通过新增 column 避免旧任务被改写；PackNet/CPG 则通过剪枝、冻结、复用和扩展来管理参数。它们可靠但资源成本较高。

% =====================================================================
\section{Memory Replay：让旧任务重新出现}
% =====================================================================

第三大方向是 replay。最朴素的 replay 是保存旧任务样本，在训练新任务时和新任务数据一起训练，近似 multi-task training。但保存原始旧数据可能很贵，也可能涉及隐私。课程重点展示了 generative replay：用生成模型产生旧任务 pseudo-data。

\framefig{0.84}{figures/fig11_generative_replay_data.jpg}
{Generative replay：为旧任务训练生成模型，用它产生 pseudo-data；学习新任务时，把生成的旧任务样本与新任务数据一起训练，近似 multi-task learning。}
{00:31:45--00:32:00}

Generative replay 的关键好处是减少原始数据存储。只要生成模型能保留旧任务分布，就可以在新任务训练时不断“回放”旧任务。但这也引入新问题：生成模型本身也可能忘记，生成质量不足会造成偏差，且生成模型的训练和存储也有成本。

\subsection{Adding New Classes：类别增量学习}

课程也提到 adding new classes 的场景。新任务不是换一个完全不同任务，而是给分类器增加新类别。例如原本只会分旧类别，现在要加入新类别；模型最终需要在所有类别之间判断，而不是只在当前任务内判断。相关方法包括 Learning without Forgetting（LwF）和 iCaRL。

\framefig{0.84}{figures/fig12_adding_new_classes_lwf_icarl.jpg}
{Adding new classes 的代表方法：LwF 使用旧模型对旧任务的响应作为目标，iCaRL 关注 class-incremental learner 和表示学习。}
{00:33:45--00:34:00}

LwF 的思想和 knowledge distillation 有关联：新任务训练时，让新模型同时匹配旧模型对旧类别的输出，从而保留旧类别知识。iCaRL 则通常维护少量 exemplar，并学习适合类别增量的 representation。Class-incremental learning 往往比 task-incremental 更难，因为测试时模型不一定知道输入属于哪个任务。

\subsection{Curriculum Learning：学习顺序也会影响遗忘}

最后课程回到学习顺序。两个任务若按 task 1 $\to$ task 2 学习，可能出现 task 1 遗忘；换成 task 2 $\to$ task 1，结果可能不同。Curriculum learning 问的是：是否存在更合适的任务顺序，让模型更少遗忘、更容易迁移？

\framefig{0.84}{figures/fig13_curriculum_learning_order.jpg}
{Curriculum learning 的问题：不同任务顺序可能导致不同遗忘与迁移结果，因此需要思考 proper learning order。}
{00:36:30--00:36:38}

这说明 continual learning 不只是算法问题，也是训练流程问题。任务顺序、任务相似度、数据平衡、replay 频率、模型容量分配，都会影响最终遗忘程度。

\subsection{本章小结}

Replay 通过让旧任务信号重新参与训练来缓解遗忘，可以保存真实样本，也可以生成 pseudo-data。新类别学习和 curriculum learning 进一步说明，continual learning 的难点不仅在参数保护，也在任务定义、类别空间和学习顺序。

% =====================================================================
\section{总结与延伸}
% =====================================================================

P54 把 catastrophic forgetting 的解决方向整理成方法地图：
\begin{itemize}
    \item \textbf{Regularization-based：} 用参数重要性保护旧任务，例如 EWC、SI、MAS。
    \item \textbf{Gradient-constrained：} 用旧任务 memory 约束当前梯度，例如 GEM。
    \item \textbf{Resource allocation：} 通过新增、冻结、剪枝、mask 或扩展为不同任务分配参数。
    \item \textbf{Replay-based：} 保存或生成旧任务数据，让训练过程接近 multi-task learning。
    \item \textbf{Curriculum-aware：} 调整任务顺序，利用任务之间的迁移关系减少遗忘。
\end{itemize}

这些方法没有一个单独完美。正则化省存储但可能保护不足；GEM 和 replay 更直接但需要旧数据或 memory；资源分配稳定但模型会增长；生成式 replay 省原始数据但依赖生成质量；curriculum 有帮助但不总能控制任务到达顺序。

\begin{importantbox}{Continual learning 的本质 trade-off}
    旧任务保护越强，新任务学习越可能受限；新任务学习越自由，旧任务越容易被破坏。所有方法都在 stability、plasticity、memory、compute、model size 与 task transfer 之间做权衡。
\end{importantbox}

从 P53 到 P54，课程完成了一个闭环：先定义灾难性遗忘和评价指标，再给出缓解方法。后续真正设计 continual learning 系统时，应先明确约束：是否允许保存旧数据、任务边界是否已知、类别空间是否增长、模型是否可扩展、旧任务表现要保护到什么程度。方法选择必须服从这些约束，而不是只看单一 benchmark 分数。

\subsection{本章小结}

P54 的核心是方法地图：保护重要参数、约束梯度、分配新资源、回放旧数据、调整学习顺序。它们从不同侧面缓解 catastrophic forgetting，但都无法绕开稳定性与可塑性的根本张力。

\end{document}
